\documentclass[11pt]{article}
\usepackage[utf8]{inputenc}
\usepackage[T1]{fontenc}
\usepackage[spanish]{babel}
\usepackage{amsmath,amssymb,amsthm}
\usepackage[margin=2.5cm]{geometry}

\newtheorem*{afirmacion}{Afirmación}
\renewcommand{\proofname}{Demostración}
\newcommand{\Z}{\mathbb{Z}}
\newcommand{\R}{\mathbb{R}}

\title{Solucionario Oficial del Grupo\\[2pt]\large Auditoría de Razonamiento Lógico en Matemática Discreta}
\date{}

\begin{document}
\maketitle

%==================================================================
\section*{Problema 1: Prueba directa y condición vacua}

\begin{afirmacion}
Para todo $n\in\Z$, si $n^2+3n+5$ es par, entonces $n$ es impar.
\end{afirmacion}

\begin{proof}
Sea $n\in\Z$ y supongamos que $n^2+3n+5$ es par (hipótesis directa).

\textbf{Paso 1: $n(n+3)$ es par para todo $n\in\Z$.}
Si $n$ es par, $n(n+3)$ es par porque tiene un factor par.
Si $n$ es impar, entonces $n+3$ es par, y $n(n+3)$ es par por la misma razón.

\textbf{Paso 2: la hipótesis es imposible.}
Por el paso 1, $n(n+3)=2m$ con $m\in\Z$. Entonces
\[
n^2+3n+5 = n(n+3)+5 = 2m+5 = 2(m+2)+1,
\]
que es impar. Esto contradice que $n^2+3n+5$ sea par, pues ningún entero es a la vez par e impar.

\textbf{Paso 3: conclusión.}
El antecedente «$n^2+3n+5$ es par» es falso para todo $n\in\Z$. Una implicación con antecedente falso es verdadera, de modo que la afirmación se cumple para todo $n$ (verdad vacua).
\end{proof}

%==================================================================
\section*{Problema 2: Contradicción y leyes de De Morgan}

\begin{afirmacion}
Si $a,b\in\Z$ y $a+b\ge 15$, entonces $a\ge 8$ o $b\ge 8$.
\end{afirmacion}

\begin{proof}
Supongamos, por contradicción, que $a+b\ge 15$ y que la conclusión es falsa. Por De Morgan,
\[
\neg(a\ge 8 \lor b\ge 8)\;\equiv\;(a<8)\land(b<8).
\]
Como $a,b$ son enteros, $a<8\Rightarrow a\le 7$ y $b<8\Rightarrow b\le 7$. Sumando,
\[
a+b\le 14,
\]
lo cual contradice $a+b\ge 15$. Por tanto $a\ge 8$ o $b\ge 8$.
\end{proof}

%==================================================================
\section*{Problema 3: Inducción matemática}

\begin{afirmacion}
Para todo $n\ge 1$: $\displaystyle\sum_{i=1}^{n} i^3=\left[\frac{n(n+1)}{2}\right]^2$.
\end{afirmacion}

\begin{proof}
Sea $P(n)$ la igualdad del enunciado.

\textbf{Caso base ($n=1$).} $\sum_{i=1}^{1} i^3=1$ y $\left[\frac{1\cdot 2}{2}\right]^2=1$. Se cumple $P(1)$.

\textbf{Paso inductivo.} Sea $k\ge 1$ y supongamos $P(k)$ (hipótesis inductiva, HI):
\[
\sum_{i=1}^{k} i^3=\left[\frac{k(k+1)}{2}\right]^2 .
\]
Separando el último sumando,
\[
\sum_{i=1}^{k+1} i^3=\underbrace{\sum_{i=1}^{k} i^3}_{\text{aquí se usa la HI}}+(k+1)^3
=\frac{k^2(k+1)^2}{4}+(k+1)^3 .
\]
Sacando factor común $\frac{(k+1)^2}{4}$:
\[
=\frac{(k+1)^2}{4}\bigl[k^2+4(k+1)\bigr]
=\frac{(k+1)^2}{4}(k^2+4k+4)
=\frac{(k+1)^2(k+2)^2}{4}
=\left[\frac{(k+1)(k+2)}{2}\right]^2,
\]
que es exactamente $P(k+1)$. Por inducción, $P(n)$ vale para todo $n\ge 1$.
\end{proof}

%==================================================================
\section*{Problema 4: Biyectividad e inversa}

Sea $f:\R\setminus\{2\}\to\R\setminus\{1\}$, $f(x)=\dfrac{x+1}{x-2}$.

\textbf{Buena definición.} Si $f(x)=1$ para algún $x\neq 2$, entonces $x+1=x-2$, es decir $1=-2$, absurdo. Luego $f(x)\neq 1$ y $f$ está bien definida en el codominio dado.

\textbf{Inyectividad.} Sean $x_1,x_2\neq 2$ con $f(x_1)=f(x_2)$. Como los denominadores son no nulos,
\[
(x_1+1)(x_2-2)=(x_2+1)(x_1-2)
\;\Longrightarrow\; -2x_1+x_2=-2x_2+x_1
\;\Longrightarrow\; 3x_2=3x_1,
\]
así que $x_1=x_2$. Luego $f$ es inyectiva.

\textbf{Sobreyectividad.} Sea $y\in\R\setminus\{1\}$. Resolvemos $\frac{x+1}{x-2}=y$:
\[
x+1=yx-2y \;\Longrightarrow\; x(y-1)=2y+1 \;\Longrightarrow\; x=\frac{2y+1}{y-1},
\]
donde se dividió entre $y-1\neq 0$. Falta comprobar que $x\neq 2$ (restricción del dominio): si
$\frac{2y+1}{y-1}=2$, entonces $2y+1=2y-2$, es decir $3=0$, absurdo. Así $x\in\R\setminus\{2\}$ y $f(x)=y$. Luego $f$ es sobreyectiva.

\textbf{Inversa.} $f$ es biyectiva, y
\[
f^{-1}:\R\setminus\{1\}\to\R\setminus\{2\},\qquad f^{-1}(y)=\frac{2y+1}{y-1}.
\]
Dominio: $\R\setminus\{1\}$. Codominio: $\R\setminus\{2\}$.

%==================================================================
\section*{Problema 5: Relación de equivalencia en $S=\{1,\dots,12\}$}

$(a,b)\in R\iff 5\mid (a^2-b^2)$.

\textbf{Reflexiva.} $a^2-a^2=0=5\cdot 0$.

\textbf{Simétrica.} Si $a^2-b^2=5k$, entonces $b^2-a^2=5(-k)$.

\textbf{Transitiva.} Si $a^2-b^2=5k$ y $b^2-c^2=5m$, sumando $a^2-c^2=5(k+m)$.

Por tanto $R$ es una relación de equivalencia. Además, $(a,b)\in R$ si y solo si $a^2\equiv b^2 \pmod 5$.

\textbf{Clases.} Calculamos $a^2 \bmod 5$:
\[
\begin{array}{c|cccccccccccc}
a & 1&2&3&4&5&6&7&8&9&10&11&12\\\hline
a^2 \bmod 5 & 1&4&4&1&0&1&4&4&1&0&1&4
\end{array}
\]
Agrupando por residuo:
\[
[5]=\{5,10\},\qquad [1]=\{1,4,6,9,11\},\qquad [2]=\{2,3,7,8,12\},
\]
con $2+5+5=12$ elementos, disjuntas y cubriendo $S$. Luego
\[
S/R=\bigl\{\{5,10\},\{1,4,6,9,11\},\{2,3,7,8,12\}\bigr\},\qquad |S/R|=3.
\]

%==================================================================
\section*{Problema Bonus: Disposiciones asturianas}

\subsection*{a) El número de disposiciones es un cuadrado perfecto}

Una disposición asturiana equivale a una función $f:S\to S$ tal que $f(s)$ es adyacente a $s$ para todo $s$ y $f$ es biyectiva (cada casilla recibe exactamente un saltamontes).

Coloreamos la cuadrícula como tablero de ajedrez y sean $B$ y $N$ las casillas blancas y negras de $S$. Casillas adyacentes tienen colores distintos, así que $f(B)\subseteq N$ y $f(N)\subseteq B$.

\textbf{Paso 1: $f|_B:B\to N$ es biyección.} Es inyectiva por serlo $f$. Es sobreyectiva porque cada $n\in N$ tiene preimagen en $S$ y esa preimagen no puede estar en $N$ (las imágenes de $N$ caen en $B$), así que está en $B$. Análogamente, $f|_N:N\to B$ es biyección.

\textbf{Paso 2: correspondencia con pares.} Sea $m$ el número de biyecciones $g:B\to N$ con $g(b)$ adyacente a $b$ para todo $b$. Recíprocamente, dado un par $(g,h)$ de biyecciones $g:B\to N$ y $h:N\to B$ que cumplen la adyacencia, la función $f=g\cup h$ es una disposición asturiana. Esta correspondencia es biyectiva.

\textbf{Paso 3: ambos factores valen $m$.} La asignación $h\mapsto h^{-1}$ lleva biyecciones $N\to B$ con adyacencia en biyecciones $B\to N$ con adyacencia (la adyacencia es simétrica), y es invertible. Luego también hay $m$ biyecciones $N\to B$.

Por el principio multiplicativo, el total es $m\cdot m=m^2$, un cuadrado perfecto. (Si $|B|\neq|N|$, $m=0$ y el total es $0=0^2$.)

\subsection*{b) La figura de 12 casillas}

\emph{Supuesto sobre la figura:} un bloque de $2\times 4$ (filas 1--2, columnas 1--4) con dos patas de dos casillas, $(3,1),(4,1)$ y $(3,4),(4,4)$. Con la paridad de $i+j$:
\[
B=\{(1,1),(1,3),(2,2),(2,4),(3,1),(4,4)\},\quad
N=\{(1,2),(1,4),(2,1),(2,3),(3,4),(4,1)\}.
\]
Contamos las biyecciones $g:B\to N$ con adyacencia:
\begin{itemize}
\item $(4,4)$ solo tiene vecino $(3,4)$, así que $g(4,4)=(3,4)$.
\item $(4,1)\in N$ solo tiene vecino $(3,1)$, así que $g(3,1)=(4,1)$.
\item Quedan $\{(1,1),(1,3),(2,2),(2,4)\}\to\{(1,2),(1,4),(2,1),(2,3)\}$ en el bloque $2\times4$.
\end{itemize}
Casos según $g(1,1)$:
\begin{enumerate}
\item $g(1,1)=(2,1)$. Si $g(2,2)=(1,2)$: $g(1,3)=(1,4),g(2,4)=(2,3)$ o $g(1,3)=(2,3),g(2,4)=(1,4)$ (2 formas). Si $g(2,2)=(2,3)$: $(1,2)$ solo puede venir de $(1,3)$, y entonces $g(2,4)=(1,4)$ (1 forma). Subtotal: $3$.
\item $g(1,1)=(1,2)$. Entonces $(2,1)$ solo puede venir de $(2,2)$. Luego $g(1,3)=(1,4),g(2,4)=(2,3)$ o $g(1,3)=(2,3),g(2,4)=(1,4)$ (2 formas). Subtotal: $2$.
\end{enumerate}
Así $m=3+2=5$ y el número de disposiciones asturianas es
\[
m^2=5^2=\boxed{25}.
\]
(Verificado además por enumeración computacional exhaustiva.)

\end{document}
